\begin{thebibliography}{99}

    \bibitem{Piezo-électricité}
        \textit{Protocole de laboratoire}, cours de matériaux, Hepia MT2 

\end{thebibliography}

\newpage

\section{Annexes}


\subsection{Hysterese et fluage}

Mesures de A, B, C, D pour le coefficient de non-linearite :


\begin{table}[h]
\centering
    \caption{Coordonnées des points A, B, C et D (arondies à la 4ème décimale) pour différentes tensions et fréquences.}
\label{tab:points}
\begin{tabular}{|c|c|c|c|c|}
\hline
\textbf{Mesure} & \textbf{Point A (x, y)} & \textbf{Point B (x, y)} & \textbf{Point C (x, y)} & \textbf{Point D (x, y)} \\
\hline
1V 5Hz & (1.0749, 0.0135) & (1.9512, 0.0359) & (1.3565, 0.0253) & (1.4771, 0.0205) \\
\hline
1V 10Hz & (1.1306, 0.0156) & (1.8983, 0.0361) & (1.1652, 0.0235) & (1.9453, 0.0320) \\
\hline
1V 20Hz & (1.2186, 0.0172) & (1.7388, 0.0324) & (0.9923, 0.0209) & (1.9767, 0.0282) \\
\hline
2V 5Hz & (0.6478, 0.0034) & (2.3305, 0.0490) & (0.5205, 0.0073) & (1.6298, 0.0190) \\
\hline
2V 10Hz & (0.8405, 0.0052) & (2.2170, 0.0415) & (0.5094, 0.0095) & (1.7278, 0.0163) \\
\hline
2V 20Hz & (0.9748, 0.0110) & (2.0367, 0.0391) & (0.5036, 0.0167) & (2.4636, 0.0324) \\
\hline
3V 5Hz & (0.3983, -0.0033) & (2.8626, 0.0531) & (1.6277, 0.0385) & (1.8169, 0.0181) \\
\hline
3V 10Hz & (0.5721, -0.0014) & (2.6232, 0.0488) & (-0.0009, 0.0043) & (2.0623, 0.0176) \\
\hline
3V 20Hz & (0.6218, 0.0040) & (2.2958, 0.0426) & (-0.0012, 0.0115) & (2.3075, 0.0176) \\
\hline
\end{tabular}
\end{table}

Calcul incertitude : 

Le coefficient de non-linéarité est défini par :
\[
K = \frac{\text{CD}}{\text{AB}}
\]

Les incertitudes sur \(\text{AB}\) et \(\text{CD}\) sont :
\[
\Delta \text{AB} = \pm 0.02 \, \text{V}, \quad \Delta \text{CD} = \pm 0.06 \, \text{V}
\]

L'incertitude relative sur \(K\) est donnée par :
\[
\frac{\Delta K}{K} = \sqrt{\left( \frac{\Delta \text{CD}}{\text{CD}} \right)^2 + \left( \frac{\Delta \text{AB}}{\text{AB}} \right)^2}
\]

L'incertitude absolue sur \(K\) est :
\[
\Delta K = K \times \sqrt{\left( \frac{\Delta \text{CD}}{\text{CD}} \right)^2 + \left( \frac{\Delta \text{AB}}{\text{AB}} \right)^2}
\]

\[
\text{AB} = 0.022\,\text{V}, \quad \text{CD} = 0.005\,\text{V}, \quad K = \frac{0.005}{0.022} \approx 0.227
\]
\[
\frac{\Delta K}{K} = \sqrt{\left( \frac{0.06}{0.005} \right)^2 + \left( \frac{0.02}{0.022} \right)^2} \approx 12.03
\]
\[
\Delta K = 0.227 \times 12.03 \approx 2.73
\]


\subsection{Interferometrie}



\begin{figure}[H]
    \centering
    \includegraphics[scale=0.5]{5V.jpg}
    \label{fig:franges}
    \caption{Interferances}
\end{figure}


\begin{figure}[H]
    \centering
    \includegraphics[scale=0.5]{10V.jpg}
    \label{fig:franges}
    \caption{Interferances}
\end{figure}


\begin{figure}[H]
    \centering
    \includegraphics[scale=0.5]{15V.jpg}
    \label{fig:franges}
    \caption{Interferances}
\end{figure}


\begin{figure}[H]
    \centering
    \includegraphics[scale=0.5]{20V.jpg}
    \label{fig:franges}
    \caption{Interferances}
\end{figure}


\begin{figure}[H]
    \centering
    \includegraphics[scale=0.5]{25V.jpg}
    \label{fig:franges}
    \caption{Interferances}
\end{figure}


\begin{figure}[H]
    \centering
    \includegraphics[scale=0.5]{30V.jpg}
    \label{fig:franges}
    \caption{Interferances}
\end{figure}


